% =================================================
% Справочный лист
% =================================================

\begin{center}
  {\fontsize{25}{30}\selectfont\bfseries\sffamily\color{Primary}
    Точки на отрезке}

  \vspace{0.3em}

  {\large\sffamily\color{GrayText}
    Порядок точек, части отрезка и несколько решений}
\end{center}

\vspace{0.45em}

\section{Какая точка лежит между двумя другими}

\betweenfigure

\begin{propertybox}[Порядок точек]
  На рисунке точка \(C\) лежит между точками \(A\) и \(B\). Порядок точек
  можно кратко записать так:
  \[
    A-C-B.
  \]
\end{propertybox}

Записи \(A-C-B\) и \(B-C-A\) описывают одно и то же расположение: в обоих
случаях точка \(C\) находится между двумя другими точками.

\begin{warningbox}[Рисунок нужно читать внимательно]
  Буквы в названии отрезка не сообщают порядок внутренних точек. Порядок
  устанавливают по рисунку или по данным задачи.
\end{warningbox}

\section{Целый отрезок и его части}

\wholepartsfigure

\begin{propertybox}[Сложение длин]
  Если точка \(C\) лежит между \(A\) и \(B\), то
  \[
    AB=AC+CB.
  \]
\end{propertybox}

Отсюда:
\[
  AC=AB-CB,\qquad CB=AB-AC.
\]

\begin{notebox}[Сначала определите целое]
  Складывают длины частей, когда находят весь отрезок. Вычитают, когда
  находят недостающую часть.
\end{notebox}

\section{Две точки внутри отрезка}

\fourpointsfigure

При порядке \(A-C-D-B\) отрезок \(AB\) состоит из трёх частей:
\[
  AB=AC+CD+DB.
\]
Поэтому
\[
  CD=AB-AC-DB.
\]

\oppositeendsfigure

\begin{propertybox}[Расстояния от разных концов]
  Если \(AC\) отложено от конца \(A\), а \(DB\) --- от конца \(B\), сначала
  сравните сумму \(AC+DB\) с длиной \(AB\).
\end{propertybox}

\threeoppositecasesfigure

\begin{center}
  \renewcommand{\arraystretch}{1.38}
  \begin{tabularx}{\textwidth}{
    >{\raggedright\arraybackslash}p{0.26\textwidth}
    >{\raggedright\arraybackslash}p{0.28\textwidth}
    >{\raggedright\arraybackslash}X}
    \toprule
    \rowcolor{TableHeader}
    Сравнение & Расположение & Расстояние между точками \\
    \midrule
    \(AC+DB<AB\)
      & между \(C\) и \(D\) есть промежуток
      & \(CD=AB-AC-DB\) \\
    \addlinespace[0.2em]
    \(AC+DB=AB\)
      & точки \(C\) и \(D\) совпадают
      & \(CD=0\) \\
    \addlinespace[0.2em]
    \(AC+DB>AB\)
      & точка \(D\) расположена левее \(C\)
      & из большей части вычитают меньшую \\
    \bottomrule
  \end{tabularx}
\end{center}

\section{Расстояния от одного конца}

\sameendpointfigure

Если известны расстояния \(AC\) и \(AD\) от одной и той же точки \(A\),
то расстояние между \(C\) и \(D\) находят вычитанием меньшей длины из
большей. На рисунке:
\[
  CD=AD-AC=9-4=5\text{ см}.
\]

\begin{warningbox}[Почему нельзя просто сложить?]
  Отрезки \(AC\) и \(AD\) имеют общую часть. Отрезок \(AC\) уже входит в
  отрезок \(AD\), поэтому для нахождения \(CD\) выполняют вычитание.
\end{warningbox}

\section{Когда задача имеет несколько решений}

\twopositionsfigure

Если известно только, что \(CD\) имеет заданную длину, точку \(D\) можно
попытаться отметить по обе стороны от \(C\). Получаются два возможных
положения.

\candidatefilterfigure

\begin{propertybox}[Алгоритм перебора]
  \begin{enumerate}
    \item Начертите отрезок \(AB\) и отметьте известную точку \(C\).
    \item От точки \(C\) отложите длину \(CD\) влево и вправо.
    \item Оставьте только положения, удовлетворяющие условию задачи.
    \item Для каждого оставшегося случая отдельно найдите ответ.
  \end{enumerate}
\end{propertybox}

Задача может иметь два решения, одно решение или не иметь решений.

\section{Краткая памятка}

\begin{center}
  \renewcommand{\arraystretch}{1.38}
  \begin{tabularx}{\textwidth}{
    >{\raggedright\arraybackslash}p{0.34\textwidth}
    >{\raggedright\arraybackslash}X}
    \toprule
    \rowcolor{TableHeader}
    Ситуация & Действие \\
    \midrule
    Одна точка между концами
      & использовать \(AB=AC+CB\) \\
    \addlinespace[0.2em]
    Две точки, расстояния от одного конца
      & из большей длины вычесть меньшую \\
    \addlinespace[0.2em]
    Две точки, расстояния от разных концов
      & сравнить \(AC+DB\) с \(AB\) \\
    \addlinespace[0.2em]
    Известно расстояние от одной внутренней точки до другой
      & проверить оба направления от исходной точки \\
    \bottomrule
  \end{tabularx}
\end{center}

\begin{questionsbox}[Проверьте себя]
  \begin{enumerate}
    \item Что означает запись \(A-C-B\)?
    \item Как найти часть отрезка, если известны весь отрезок и другая часть?
    \item Как найти расстояние между двумя точками, если известны их расстояния от \(A\)?
    \item Почему в некоторых задачах нужно рассматривать два рисунка?
    \item Как проверить, подходит ли найденное положение точки условию задачи?
  \end{enumerate}
\end{questionsbox}
